% !TeX root = ../sections/02_exercises.tex
\subsection{2-B-7.}
\begin{exercise}
	実数
	\[
	y_1>x_1>0
	\]
	に対して，数列
	\[
	\{x_n\}_{n\in\mathbb{N}},
	\qquad
	\{y_n\}_{n\in\mathbb{N}}
	\]
	を以下の漸化式を用いて定義する。
	\[
	x_{n+1}:=\sqrt{x_ny_n},
	\qquad
	y_{n+1}:=\frac{x_n+y_n}{2}
	\qquad
	(n=1,2,3,\ldots).
	\]
	
	以下の問いに答えよ。
	
	\begin{enumerate}
		\item
		任意の \(n\in\mathbb{N}\) に対して，
		\[
		x_n\leq x_{n+1}
		\leq y_{n+1}
		\leq y_n
		\]
		が成立することを証明せよ。
		
		\item
		数列 \(\{x_n\}_{n\in\mathbb{N}}\)，
		\(\{y_n\}_{n\in\mathbb{N}}\) は共に収束し，
		それらの極限値は一致することを証明せよ。
	\end{enumerate}
\end{exercise}

\begin{background}
	この問題では，相加平均と相乗平均の関係を用いる。
	正の実数 \(x,y\) に対して，
	\[
	\sqrt{xy}\leq\frac{x+y}{2}
	\]
	が成り立つ。
	これは AM-GM 不等式である。
	
	また，\(0<x\leq y\) のとき，
	\[
	x\leq\sqrt{xy}\leq y
	\]
	も成り立つ。
	
	したがって，
	\[
	x_{n+1}=\sqrt{x_ny_n},
	\qquad
	y_{n+1}=\frac{x_n+y_n}{2}
	\]
	は，\(x_n\) と \(y_n\) の間にある二つの平均を作る操作である。
	
	この操作を繰り返すと，
	\(\{x_n\}\) は上に上がり，
	\(\{y_n\}\) は下に下がる。
	したがって，単調収束定理を用いて収束性を示すことができる。
\end{background}

\begin{strategy}
	まず，帰納法により
	\[
	0<x_n\leq y_n
	\]
	を確認する。
	初期条件より
	\[
	0<x_1<y_1
	\]
	である。
	
	いま \(0<x_n\leq y_n\) とする。
	このとき
	\[
	x_n^2\leq x_ny_n
	\]
	であり，両辺の平方根を取ると
	\[
	x_n\leq \sqrt{x_ny_n}=x_{n+1}
	\]
	となる。
	
	また，AM-GM 不等式より
	\[
	x_{n+1}
	=
	\sqrt{x_ny_n}
	\leq
	\frac{x_n+y_n}{2}
	=
	y_{n+1}
	\]
	である。
	
	さらに \(x_n\leq y_n\) より
	\[
	\frac{x_n+y_n}{2}\leq y_n
	\]
	であるから，
	\[
	y_{n+1}\leq y_n
	\]
	となる。
	
	以上より
	\[
	x_n\leq x_{n+1}
	\leq y_{n+1}
	\leq y_n
	\]
	が成り立つ。
	
	次に，収束性を示す。
	上の不等式から，\(\{x_n\}\) は単調増加であり，
	\[
	x_n\leq y_n\leq y_1
	\]
	なので上に有界である。
	したがって，\(\{x_n\}\) は収束する。
	
	また，\(\{y_n\}\) は単調減少であり，
	\[
	x_1\leq x_n\leq y_n
	\]
	なので下に有界である。
	したがって，\(\{y_n\}\) も収束する。
	
	そこで
	\[
	\lim_{n\to\infty}x_n=\alpha,
	\qquad
	\lim_{n\to\infty}y_n=\beta
	\]
	とおく。
	
	漸化式
	\[
	y_{n+1}=\frac{x_n+y_n}{2}
	\]
	の両辺で \(n\to\infty\) とすると，
	\[
	\beta=\frac{\alpha+\beta}{2}
	\]
	となる。
	したがって
	\[
	2\beta=\alpha+\beta
	\]
	より
	\[
	\alpha=\beta
	\]
	が従う。
	
	よって，二つの数列は共に収束し，極限値は一致する。
\end{strategy}